\scanpage{117}
\mspage{32}
% archive sheet number: 116

% L01-05
et l'image inverse $U(m)$ de $U$ dans $X(m)$, qui est un rev. étale \deleted{\unsure{pr.}} principal de $U$ de groupe $\mu_m(X)^r=G(m)$, les $U(m)$ forment un syst. projectif de \deleted{\unsure{schémas affines}} \unsure{revêts} étales \deleted{\unsure{et}} principaux, compatibles avec les hom de transition évidents

% L06-07
\diagram{\begin{tikzcd}[ampersand replacement=\&]
\mu_{m'}^r(X)\arrow[r,no head]\arrow[d,equal]\&\mu_m^r(X)\arrow[d,equal]\\
G(m')\&G(m)
\end{tikzcd}}
\inserted{pour} $(m\mid m')$.

% L08-09
Soit \deleted{$\unsure{U(\infty)}$}
\[
U(\infty)=\varprojlim U(m)
\]

% L10-11
qui est un \deleted{\unsure{revêtement proétale de}} pro-revêtement étale de $U$ de groupe

% L12
\[
G\deleted{(\unsure{\infty})}=\varprojlim_m G(m)=\prod_{\ell\ne p}\unsure{\mathbb{Z}|Z}_\ell(1)^r,
\tag*{(4.8.1)}
\]

% L13-14
\inserted{$p$ est la car. résiduelle de $X$, et où} où l'on pose par définition

% L15
\[
\mathbb{Z}_\ell(1)=\varprojlim_\nu\mu_{\ell^\nu}(X).
\tag*{(4.8.2)}
\]

% L16-17
Soit $n$ un entier $>0$ fixé tel que $n$ soit premier à $p$, \inserted{et soit encore $\Lambda=\mathbb{Z}/n\mathbb{Z}$.} On a donc

% L18
\[
H^*(U(\infty),\Lambda)\simeq\varinjlim_m H^*(U(m),\Lambda),
\tag*{\unsure{(*)|(1)}}
\]

% L19-23
\deleted{et \unsure{compte tenu} que}
\[
\deleted{U(m)=X(m)-\bigcup_{1\leq i\leq r}D'_i(m),}
\]
\deleted{avec $X(m)$ régulier et les $D'_i(m)\inserted{=\operatorname{div}(T_i(m))}$ des diviseurs réguliers dont la somme est à croisements normaux, on trouve que}

% L23-24
D'autre part, d'après les rappels de 4.6, le calcul de l'hom de transition
